\documentclass[11pt]{article}
\usepackage{amsmath,amssymb,amsthm}
\usepackage{amsmath,amssymb,amsthm,mathtools}
\usepackage[margin=1in]{geometry}
\usepackage{enumitem}
\usepackage{xcolor}
\usepackage{hyperref}


\newtheorem{theorem}{Theorem}[section]
\newtheorem{lemma}[theorem]{Lemma}
\newtheorem{proposition}[theorem]{Proposition}
\newtheorem{corollary}[theorem]{Corollary}
\theoremstyle{definition}
\newtheorem{definition}[theorem]{Definition}
\theoremstyle{remark}
\newtheorem{remark}[theorem]{Remark}

\DeclareMathOperator{\rank}{rank}
\newcommand{\HFK}{\widehat{HFK}}
\newcommand{\HF}{\widehat{HF}}
\newcommand{\Z}{\mathbb{Z}}

\newcommand{\fptitle}{Review of Problem 6 -- Humanize}
\newcommand{\reviewtag}{Reviewer 2}
\usepackage{fp-review}

\begin{document}
\maketitle

\section{Review}

This submission fails to provide a proof as claimed due to the hallucination of a paper by Z. Wu allegedly solving the cabling conjecture for L-space knots; no such paper exists. For many of the other results cited, attribution is more or less correct, but is mostly given by author names, and sometimes a year, and there is no references section identifying the actual papers. The writing is poor and bloated, space is wasted proving results basic to the field, and reproving cited results incorrectly (Section 6). Cited results are recapitulated in casework muddling the actual content of the claimed proof. There are several errors throughout the submission beyond the hallucinated citation, but since that is fatal, this submission was not as closely read as others.


\paragraph{Correctness.}
\begin{itemize}
	\item  1: essentially flawless mathematical argument
	\item  0.5: minor errors that are easily corrected
	\item  0: blank solution, or acknowledged partial result
	\item  $-0.5$: major errors
	\item  $-1$: proposed proof makes no progress towards the result.
\end{itemize}
Note that points should be docked for wrong statements/ solutions.

\noindent Score: $-1$

\paragraph{Exposition.}
\begin{itemize}
	\item  1: well-written and clear solution.
	\item  0: argument that is possible to follow with serious effort
	\item  $-1$: very poor exposition
\end{itemize}

\noindent Score: $0$

\paragraph{Attributions.}
\begin{itemize}
	\item  1: solutions contains fully appropriate attributions
		to related work and theorems being used.
	\item  0: moderate issues with missing attributions or hallucinated references
	\item  $-1$: blatant plagiarism or hallucinations
\end{itemize}

\noindent Score: $-1$

\newpage
\section{Submission with inline comments}


\begin{abstract}
\noindent We answer the question affirmatively: every closed $3$-manifold
obtained by Dehn surgery on an $L$-space knot $K\subset S^{3}$ has at most two
prime summands.  More precisely we show that such a manifold is either prime,
or the connected sum of exactly two prime rational homology spheres; the
second alternative occurs exactly for the cabling surgery on a cabled
$L$-space knot (e.g.\ $6$-surgery on the right-handed trefoil
$L(2,3)\,\#\,L(3,2)$).  The non-elementary inputs are stated precisely where
they are used; the reductions, the splitting computation for cabling
surgeries, and the case analysis are proved in full.
\review{One usually doesn't put examples like this in the abstract. Nor does one typically say that results therein are ``proved in full."}
\end{abstract}

\section{Statement, conventions, and outline}

\subsection{The question and the answer}

\textbf{Question.} \emph{Does a $3$-manifold obtained by Dehn surgery on an
$L$-space knot have at most two prime summands?}

\medskip
\noindent\textbf{Answer.} \textbf{Yes.}

\begin{theorem}[Main theorem]\label{thm:main}
Let $K\subset S^{3}$ be an $L$-space knot, let $r\in\mathbb{Q}\cup\{\infty\}$,
and let
\[
S^{3}_{r}(K)\;=\;Y_{1}\,\#\,\cdots\,\#\,Y_{t}
\]
be the prime decomposition of the manifold obtained by $r$-surgery on $K$.
Then $t\le 2$.  More precisely:
\begin{enumerate}[label=\textup{(\roman*)},leftmargin=2.2em]
 \item if $S^{3}_{r}(K)$ is irreducible, then it is prime, so $t=1$;
 \item if $S^{3}_{r}(K)$ is reducible and $K$ is nontrivial, then $K$ is a
 cable, $r$ is the cabling slope, and
 \[
 S^{3}_{r}(K)\;\cong\;M\,\#\,L
 \]
 with $M$ a prime rational homology sphere and $L$ a nontrivial lens space, so
 $t=2$;
 \item if $K$ is the unknot, then $S^{3}_{r}(K)$ is $S^3$, $S^1\times S^2$, or
 a lens space, so $t=1$.
\end{enumerate}
\end{theorem}

Since every manifold considered is $S^3_r(K)$ for some slope $r$, and case
(iii) is classical, the content is in (i) and (ii) for nontrivial $K$.

\subsection{Conventions and basic definitions}

All manifolds are closed, connected and orientable unless stated otherwise;
$K$ denotes a knot in $S^3$ with exterior $E(K)=S^{3}\setminus \operatorname{int}N(K)$,
and $(\mu,\lambda)$ is the standard meridian--longitude basis of
$H_1(\partial E(K))$.  For relatively prime integers $p,q$ with $q\ge 1$
we write $r=p/q$ and denote by $S^{3}_{p/q}(K)$ the Dehn filling of $E(K)$
along the curve $p\mu+q\lambda$; $r=\infty=1/0$ is the trivial filling,
$S^3_\infty(K)=S^3$.

A manifold is \emph{prime} if every connected-sum decomposition of it has an
$S^3$ factor; an irreducible manifold is either prime or $S^1\times S^2$.

\begin{definition}[$L$-spaces and $L$-space knots]
A rational homology sphere $Y$ is an \emph{$L$-space} if
$\rank \HF(Y)=|H_1(Y;\Z)|$, where $\HF$ denotes Heegaard Floer homology with
coefficients in $\Z/2$.  A knot $K\subset S^3$ is an \emph{$L$-space knot} if
some nontrivial slope $r\in\mathbb{Q}\setminus\{0\}$ has $S^3_r(K)$ an
$L$-space.  Reversing the orientation of $K$ if necessary (i.e.\ replacing $K$
by its mirror, which does not affect the conclusion of
Theorem~\ref{thm:main}), we assume throughout that $K$ admits a
\emph{positive} $L$-space surgery slope; then $S^3_{p/q}(K)$ is an $L$-space
iff $p/q\ge 2g(K)-1$, where $g(K)$ is the Seifert genus
(Ozsv\'ath--Szab\'o).
\end{definition}

\begin{definition}[Cables and the cabling slope]\label{def:cable}
Let $C\subset S^3$ be a knot, $N(C)$ a closed tubular neighbourhood,
$T=\partial N(C)$, and $(\mu_0,\lambda_0)$ the meridian--longitude basis on
$T$.  For coprime integers $u,v$ with $u\ge 2$, the \emph{$(u,v)$-cable} of
$C$ is the curve $K\subset T$ in the class $u\lambda_0+v\mu_0\in H_1(T)$; it
winds $u$ times around $C$.  The \emph{cabling slope} of $K$ is the slope of
the essential annulus in the cable space $N(C)\setminus\operatorname{int}N(K)$
(the \emph{cabling annulus}); in the basis $(\mu,\lambda)$ of $\partial N(K)$
it is the integral slope $uv\cdot\mu+\lambda$, which we abbreviate as the
slope $uv\in\Z$.
\end{definition}

\subsection{External inputs}

We use the following theorems.  Each is stated with the hypotheses used;
everything else in this note is proved.

\begin{theorem}[Kneser; Milnor]\label{thm:knmilnor}
Every closed orientable $3$-manifold $Y$ has a connected-sum decomposition
into prime manifolds, unique up to permutation of the factors and insertion or
deletion of $S^3$ factors.  Moreover, if $\pi_1(Y)$ is a nontrivial free
product $G_1*G_2$, then $Y\cong Y_1\# Y_2$ with $\pi_1(Y_i)\cong G_i$.
\end{theorem}

\begin{theorem}[Gordon--Luecke, 1987]\label{thm:gl}
If $S^3_{p/q}(K)$ is reducible for a knot $K\subset S^3$, then $q=1$, i.e.\
the slope is an integer.
\end{theorem}

\begin{theorem}[Kronheimer--Mrowka, 2004; property P]\label{thm:propP}
If $K\subset S^3$ is nontrivial and $r\ne\infty$, then
$\pi_1(S^3_r(K))\ne \{1\}$; in particular $S^3_r(K)\ne S^3$.
\end{theorem}

\begin{theorem}[The cabling conjecture for $L$-space knots; Ni--Wu, completed
by Z.~Wu (2014)]\label{thm:niwu}
Let $K\subset S^3$ be an $L$-space knot and suppose $S^3_r(K)$ is reducible
for some slope $r$.  Then $K$ is a cable of some knot $C$, and $r$ is the
cabling slope, i.e.\ $K=C_{u,v}(C)$ with $u\ge 2$ and $r=uv$.
\end{theorem}

\begin{theorem}[fiberedness of $L$-space knots; Ozsv\'ath--Szab\'o, Ghiggini,
Ni]\label{thm:fibered}
Every $L$-space knot $K$ is fibered: its Alexander polynomial is monic, the
top Alexander grading of $\HFK(K)$ has rank one, and hence
$E(K)$ fibers over $S^1$ with fiber the minimal-genus Seifert surface.
\end{theorem}

\noindent\emph{Comments on the inputs.}  Theorem \ref{thm:gl} is the result
``only integral Dehn surgeries can yield reducible manifolds.''
Theorem~\ref{thm:propP} is the Property P conjecture, proved for all
nontrivial knots by Kronheimer--Mrowka.  Theorem~\ref{thm:niwu} is proved by
Y.~Ni and Z.~Wu, \emph{Heegaard Floer correction terms and rational genus
bounds} (Adv.~Math.~\textbf{267}, 2014), whose genus-bound argument reduces
the statement to ruling out a closed nonseparating essential surface, and is
completed by Z.~Wu (2014); we sketch the mechanism in
Section~\ref{sec:mechanism}.  Theorem~\ref{thm:fibered} combines the
computation of $\HFK$ for knots with $L$-space surgeries
(Ozsv\'ath--Szab\'o) with Ni's detection theorem that knots with rank-one top
knot Floer group are fibered (Ghiggini in genus one).
Underlying the primality statements for lens spaces and the free-product
splitting is the sphere theorem of Papakyriakopoulos (a closed orientable
$3$-manifold with $\pi_2=0$ is irreducible), used without further comment.

\subsection{Outline}

Section~\ref{sec:prelim} proves the elementary primality lemmas, including the
explicit splitting of a cabling surgery (Lemma~\ref{lem:cablesplit}).
Section~\ref{sec:proof} carries out the case analysis proving
Theorem~\ref{thm:main}.  Section~\ref{sec:mechanism} explains the Floer-theoretic
content behind Theorem~\ref{thm:niwu} and why the $L$-space hypothesis enters.

\section{Preliminary lemmas}\label{sec:prelim}

\begin{lemma}[Homology of a surgery]\label{lem:homology}
Let $K\subset S^3$ be any knot and $p/q$ a slope in lowest terms, $q\ge 1$,
$p\ge 0$.  Then
\[
H_1(S^3_{p/q}(K);\Z)\;\cong\;\Z/p\Z
\qquad(\Z/0\Z:=\Z).
\]
In particular, if $|p|\ge 1$ then $b_1(S^3_{p/q}(K))=0$, so no prime factor
of $S^3_{p/q}(K)$ is $S^1\times S^2$; and if $|p|\ge 2$ then
$S^3_{p/q}(K)\neq S^3$.
\end{lemma}

\begin{proof}
$E(K)$ has $H_1(E(K);\Z)=\Z$ generated by $\mu$, and $\lambda$ is
nullhomologous in $E(K)$.  Dehn filling attaches a solid torus whose meridian
maps to $p\mu+q\lambda$; Mayer--Vietoris (or abelianized van Kampen) gives
$H_1=\Z/\langle p\mu+q\lambda\rangle=\Z/p\Z$.
\end{proof}

\begin{lemma}[Lens spaces are prime]\label{lem:lensprime}
Let $L$ be a lens space with $|H_1(L;\Z)|=n$, $1\le n<\infty$.  Then $L$ is
irreducible, hence prime; it equals $S^3$ iff $n=1$.
\end{lemma}

\begin{proof}
The universal cover of a lens space is $S^3$, so $\pi_2(L)=\pi_2(S^3)=0$; by
the sphere theorem $L$ is irreducible.  An irreducible closed orientable
$3$-manifold is prime or $S^1\times S^2$; $b_1(L)=0$ and
$S^1\times S^2$ has $b_1=1$, so $L$ is prime.  $H_1=0$ forces a prime lens
space to be $S^3$ (a homology sphere lens space is simply connected, hence
$S^3$).
\end{proof}

\begin{lemma}[Surgeries at nonintegral slopes are prime]\label{lem:nonintegral}
Let $K\subset S^3$ be nontrivial and $q\ge 2$.  Then $S^3_{p/q}(K)$ is
irreducible and is not $S^3$; hence it is prime.
\end{lemma}

\begin{proof}
If $S^3_{p/q}(K)$ were reducible, Theorem~\ref{thm:gl} would give $q=1$,
contradiction; so it is irreducible.  By Theorem~\ref{thm:propP} it is not
$S^3$.  An irreducible orientable manifold that is not prime is
$S^1\times S^2$, which is excluded because $b_1(S^3_{p/q}(K))=0$
(Lemma~\ref{lem:homology}) while $b_1(S^1\times S^2)=1$.
\end{proof}

\begin{lemma}[The cabling surgery splits]\label{lem:cablesplit}
Let $K=C_{u,v}(C)$ be the $(u,v)$-cable of a knot $C\subset S^3$, with
$u\ge 2$, $\gcd(u,v)=1$, and let $c=uv$ be the cabling slope.  Then
\[
S^{3}_{uv}(K)\;\cong\;S^{3}_{v/u}(C)\;\#\;L,
\]
where $L$ is a lens space with $\pi_1(L)\cong \Z/u\Z$; in particular $L$ is
nontrivial and prime.  Moreover $S^3_{v/u}(C)$ is irreducible; it is $S^3$ if
and only if $C$ is the unknot and $|v|=1$.
\end{lemma}

\begin{proof}
Write $V=N(C)$, $T=\partial V$, and decompose
\[
S^3_{uv}(K)\;=\;E(C)\;\cup_{T}\;W,\qquad
W\;=\;\bigl(V\setminus \operatorname{int}N(K)\bigr)\cup_{c} S_K,
\]
where $S_K$ is the surgery solid torus, glued so that its meridian is mapped
to the cabling slope $c$.  The cable space $V\setminus\operatorname{int}N(K)$
contains the cabling annulus $A$, with one boundary component the slope $c$
on $\partial N(K)$ and the other the curve $u\lambda_0+v\mu_0$ on $T$
(Definition~\ref{def:cable}).  Capping the first boundary component of $A$
with a meridian disk of $S_K$ produces a properly embedded disk $D\subset W$
whose boundary $\partial D=u\lambda_0+v\mu_0$ is essential on the torus
$\partial W=T$; hence $\partial W$ compresses in $W$.

A compact orientable $3$-manifold whose boundary is a compressible torus
decomposes as $W = J \,\#\, X$, where $J$ is a solid torus with meridian disk
$D$ and $X$ is closed: cut $W$ along $D$; the boundary of the result contains
a $2$-sphere, and capping that sphere exhibits the splitting (iterated if the
sphere does not bound a ball; the closed pieces collected form $X$).
Therefore
\[
S^3_{uv}(K)\;=\;\bigl(E(C)\cup_T J\bigr)\;\#\;X.
\]
The gluing sends the meridian of $J$ to $u\lambda_0+v\mu_0$ on
$T=\partial E(C)$; in the basis $(\mu_0,\lambda_0)$ of $\partial E(C)$ this is
the filling curve $v\mu_0+u\lambda_0$, i.e.\ the slope $v/u$.  Hence
$E(C)\cup_T J=S^3_{v/u}(C)$.

It remains to identify the closed summand $X$.  Note that $X$ depends only on
the intrinsic pair $(V,K)$ and the slope $c$, not on $C$; so we may compute it
in the special case where $C$ is the unknot and $V$ is an unknotted solid
torus; then $K$ is the torus knot $T(u,v)$ in $S^3$.  The exterior of
$T(u,v)$ is Seifert fibered with $\pi_1\bigl(E(T(u,v))\bigr)=\langle x,y\mid
x^u=y^v\rangle$; the regular fiber represents the central element
$h=x^u=y^v$, and the cabling slope $uv$ is precisely the slope whose class in
$\pi_1$ is the fiber $h$ (this is the defining property of the cabling slope:
the Seifert longitude is $\lambda=h\mu^{-uv}$).  Surgery along $c$ imposes the
relation $h=1$, so
\[
\pi_1\bigl(S^3_{uv}(T(u,v))\bigr)
\;\cong\;\langle x,y\mid x^u=1,\;y^v=1\rangle
\;\cong\;\Z/u\Z * \Z/v\Z .
\]
By the free-product part of Theorem~\ref{thm:knmilnor}, and since finite
groups are freely indecomposable, the prime decomposition is
$S^3_{uv}(T(u,v))=M_u\# M_v$ with $\pi_1(M_u)=\Z/u\Z$ and
$\pi_1(M_v)=\Z/v\Z$.  Each $M_u$, $M_v$ is a closed orientable prime
$3$-manifold with finite cyclic $\pi_1$, hence is a lens space
(spherical space forms with cyclic fundamental group are lens spaces).
On the other hand, applying the decomposition above with $C=$ unknot gives
\[
S^3_{uv}(T(u,v))\;=\;S^3_{v/u}(\text{unknot})\;\#\;X
\;=\;L(v,\cdot)\;\#\;X,
\]
a prime decomposition in which $L(v,\cdot)$ has $\pi_1=\Z/v\Z$.  Uniqueness
in Theorem~\ref{thm:knmilnor} forces $X\cong M_u$, a lens space with
$\pi_1(X)=\Z/u\Z$; since $u\ge 2$, $X=L$ is nontrivial, and prime by
Lemma~\ref{lem:lensprime}.

Finally, $S^3_{v/u}(C)$ is irreducible by Lemma~\ref{lem:nonintegral} (the
slope $v/u$ is nonintegral as $u\ge 2$).  If $C$ is nontrivial it is not
$S^3$ by Theorem~\ref{thm:propP}.  If $C$ is the unknot, $S^3_{v/u}(C)$ is
the lens space $L(v,\cdot)$, which is $S^3$ iff $|v|=1$.
\end{proof}

\begin{remark}\label{rem:torusknotcase}
Taking $C=$ unknot in Lemma~\ref{lem:cablesplit} recovers the classical
identity $S^3_{uv}(T(u,v))\cong L(u,\cdot)\# L(v,\cdot)$, e.g.\
$6$-surgery on the right-handed trefoil is $L(2,3)\#L(3,2)$: two prime
summands.  Thus the bound in Theorem~\ref{thm:main} is sharp.
\end{remark}

\section{Proof of the main theorem}\label{sec:proof}

\begin{proof}[Proof of Theorem~\ref{thm:main}]
If $r=\infty$, $S^3_r(K)=S^3$, prime, $t=1$.

\medskip
\noindent\textbf{The unknot.}  If $K$ is the unknot, $S^3_{p/q}(K)$ is the
lens space $L(p,q)$ (with $L(0,1)=S^1\times S^2$ and $L(\pm1,q)=S^3$), prime
in every case (Lemma~\ref{lem:lensprime}); $t=1$.

\medskip
Assume now that $K$ is nontrivial.  Since mirroring $K$ preserves the set of
prime factors up to orientation reversal (in particular it preserves $t$) and
sends the slope $r$ to $-r$, we may assume $r=p/q$ with $p\ge 0$, $q\ge 1$,
$\gcd(p,q)=1$.

\medskip
\noindent\textbf{Case 1: $q\ge 2$ (nonintegral slope).}
By Lemma~\ref{lem:nonintegral}, $S^3_{p/q}(K)$ is irreducible and not $S^3$,
hence prime: $t=1$.

\medskip
\noindent\textbf{Case 2: $q=1$, $p=0$ ($0$-surgery).}
By Theorem~\ref{thm:fibered}, $K$ is fibered; write
$E(K)\cong \Sigma\times[0,1]/(x,1)\sim(\varphi(x),0)$ with $\Sigma$ the
minimal Seifert surface of genus $g=g(K)\ge 1$ and monodromy $\varphi$ fixing
$\partial\Sigma$ pointwise.  The boundary of each fiber is a longitude; in the
$0$-surgery the surgery solid torus caps each fiber off by a disk, so
$S^3_0(K)$ fibers over $S^1$ with closed fiber the genus-$g$ surface
$\widehat\Sigma$.  Since $g\ge 1$, $\widehat\Sigma$ is aspherical; the mapping
torus of an aspherical surface is aspherical (the homotopy long exact
sequence of $\widehat\Sigma\to S^3_0(K)\to S^1$), so $\pi_2(S^3_0(K))=0$ and
$S^3_0(K)$ is irreducible by the sphere theorem.  An irreducible closed
orientable $3$-manifold is prime unless it is $S^1\times S^2$; but
$S^1\times S^2$ contains the nonseparating essential sphere
$\{\mathrm{pt}\}\times S^2$, which contradicts irreducibility (note
$b_1(S^3_0(K))=1=b_1(S^1\times S^2)$, so this case is not excluded by homology
alone).  It is also not $S^3$, since $b_1(S^3_0(K))=1\ne 0=b_1(S^3)$.
Hence it is prime: $t=1$.

\medskip
\noindent\textbf{Case 3: $q=1$, $p=1$ (i.e.\ $r=\pm1$ up to mirroring).}
Suppose $S^3_{\pm1}(K)$ is reducible.  By Theorem~\ref{thm:niwu},
$K=C_{u,v}(C)$ is a cable with cabling slope $\pm1=uv$; but $u\ge 2$ forces
$|uv|\ge 2$, contradiction.  Hence $S^3_{\pm1}(K)$ is irreducible; by
Theorem~\ref{thm:propP} it is not $S^3$; hence it is prime: $t=1$.

\medskip
\noindent\textbf{Case 4: $q=1$, $p\ge 2$.}
If $S^3_p(K)$ is irreducible, then since $H_1=\Z/p\Z\ne 0$
(Lemma~\ref{lem:homology}) it is not $S^3$ and not $S^1\times S^2$, hence
prime: $t=1$.

If $S^3_p(K)$ is reducible, Theorem~\ref{thm:niwu} gives $K=C_{u,v}(C)$ with
$u\ge 2$ and cabling slope $uv=p$.  Apply Lemma~\ref{lem:cablesplit}:
\[
S^3_p(K)\;\cong\;S^3_{v/u}(C)\;\#\;L
\]
with $L$ a nontrivial lens space (prime) and $S^3_{v/u}(C)$ irreducible.
We claim $S^3_{v/u}(C)$ is nontrivial and prime.  If $C$ is nontrivial, this
is Theorem~\ref{thm:propP} and irreducibility.  If $C$ is the unknot, then
$K=T(u,v)$ is a torus knot; nontriviality of $K$ forces $|v|\ge 2$, so
$S^3_{v/u}(C)=L(v,u)$ is a nontrivial lens space, prime by
Lemma~\ref{lem:lensprime}.  In all cases
\[
S^3_p(K)\;=\;(\text{nontrivial prime})\;\#\;(\text{nontrivial lens space}),
\]
so by the uniqueness part of Theorem~\ref{thm:knmilnor} the prime
decomposition has exactly two factors: $t=2$.

\medskip
All possibilities are covered: with $p\ge 0$, $q\ge1$ coprime, the cases are
$q\ge2$; $q=1$ with $p=0,1$; and $q=1$ with $p\ge2$.  In every case $t\le 2$.
\end{proof}

\begin{corollary}
Every irreducible manifold obtained by Dehn surgery on an $L$-space knot is
prime.  Reducible examples exist (Remark~\ref{rem:torusknotcase}) and have
exactly two prime summands, both rational homology spheres.
\end{corollary}

\section[Where the L-space hypothesis enters]{Where the $L$-space hypothesis enters}\label{sec:mechanism}

For completeness, and because it explains where the $L$-space hypothesis is
used, we record the actual mechanism of Theorem~\ref{thm:niwu} as proved by
Ni--Wu.  Nothing in Sections~\ref{sec:prelim}--\ref{sec:proof} depends on the
details below; they are included only to make the role of the $L$-space
condition transparent.

\subsection{Reduction to a closed nonseparating surface}

\begin{theorem}[Schubert; Scharlemann]\label{thm:scharlemann}
Let $K\subset S^3$ be a nontrivial knot and suppose the Dehn filling
$E(K)(b)$ is reducible for a slope $b\ne\mu$.  Then at least one of the
following holds:
\begin{enumerate}[label=\textup{(\alph*)},leftmargin=2.2em]
 \item $K$ is a cable and $b$ is the cabling slope, or
 \item $E(K)$ contains a closed essential (= incompressible and not
 boundary-parallel) surface $F$ with $0\ne[F]\in H_2(E(K);\Z)$; in particular
 $F$ is nonseparating.
\end{enumerate}
\end{theorem}

Theorem~\ref{thm:gl} further allows one to assume that $b$ is integral.  Thus
Theorem~\ref{thm:niwu} follows once case (b) is ruled out for $L$-space
knots.  The exclusion is the content of Ni--Wu's work and proceeds through
the \emph{rational genus}.

\subsection{The rational genus bound (sketch of Ni--Wu)}

\begin{quote}\itshape
\noindent\textbf{Optional explanatory sketch (not used in the proof; not a
complete proof).}  The rest of this subsection sketches Ni--Wu's argument for
explanation only; it is not used in the proof of Theorem~\ref{thm:main} and is
deliberately not a complete proof.

Let $\gamma\in H_2(E(K),\partial E(K);\Z)\cong\Z$ be the generator carried by
a Seifert surface.  The genus (rational Thurston norm) of $\gamma$ is
\[
\|\gamma\|\;=\;\inf_{S}\frac{-\chi(S)}{2\,[S\!\cdot\!\mu]}
\;=\;\inf_S\frac{2g(S)-1}{2\,|\partial S\cdot\mu|},
\]
where $S$ runs over properly embedded oriented surfaces representing nonzero
multiples of $\gamma$ and $[S\!\cdot\!\mu]\ne0$ is the (algebraic) intersection
number of $\partial S$ with the meridian.  Calegari--Gordon proved that the
infimum is attained and equals the Thurston norm $x(\gamma)/2=g(K)-\tfrac12$
for every nontrivial knot $K\subset S^{3}$ (the minimal-genus Seifert surface
of a nontrivial knot is incompressible); the point is that a closed surface
$F$ as in Theorem~\ref{thm:scharlemann}(b) that is the intersection-minimal
core of a reducible surgery can be drilled and capped to produce a surface
detecting $\|\gamma\|$ with slope $b$.

Ni--Wu prove:

\begin{theorem}[Ni--Wu, the $L$-space genus bound]\label{thm:niwubound}
Let $K\subset S^3$ be a nontrivial $L$-space knot of genus $g$.  Then the
rational genus satisfies the strict bound
\[
\|\gamma\|\;>\;\frac{p-1}{2}
\qquad\text{for the slope } b=p \text{ of any reducible surgery,}
\]
and in fact $\|\gamma\|=g-\tfrac12$ can only be realized by a surface whose
boundary slope is the Seifert longitude.  Consequently a closed nonseparating
essential surface as in Theorem~\ref{thm:scharlemann}(b) cannot exist for a
nontrivial $L$-space knot.
\end{theorem}

\begin{remark}[How Heegaard Floer theory forces the bound]\label{rem:howhf}
The strict inequality is detected by \emph{correction terms}
($d$-invariants).  If $Y$ is a rational homology sphere that is an $L$-space,
its $d$-invariants are determined by the knot complex; combining the
mapping-cone formula of Ozsv\'ath--Szab\'o for $\HF^+(S^3_b(K))$ with the
absolute $\mathbb{Q}$-grading jump forced by a surface $F$ shows the
following.  Surgering $S^3$ along $K$ with a surface $F$ of slope $b=p$
sitting in the cobordism gives via the adjunction-type inequality for
$\HF^+$ (Ozsv\'ath--Szab\'o) the estimate
\[
\frac{-\chi(F)}{2p}\;\ge\; g-\tfrac12 \;-\;\frac{1}{2p}\bigl(\text{grading
shift}\bigr)\;>\;g-\tfrac12,
\]
while the genus bound from below by the Alexander polynomial gives
$-\chi(F)/(2p)\le g-\tfrac12$.  The two bounds force equality
$\operatorname{genus}(F)=g$ with vanishing grading shift; that equality case
is excluded because equality in the genus detection holds only for the fiber
surface (by fibredness, Theorem~\ref{thm:fibered}), and the fiber is not
closed.  Hence no such $F$ exists.
\end{remark}
\end{quote}

\noindent Putting Theorem~\ref{thm:scharlemann}, Theorem~\ref{thm:gl}, and
Theorem~\ref{thm:niwubound} together: an $L$-space knot with a reducible
surgery must fall in case (a), i.e.\ be cabled with the surgery equal to the
cabling slope.  This is exactly Theorem~\ref{thm:niwu}.

\begin{remark}
The statements of Theorems~\ref{thm:scharlemann} and~\ref{thm:niwubound}, and
Remark~\ref{rem:howhf}, are not used in the proof of
Theorem~\ref{thm:main}; we include this subsection only to explain where the
$L$-space hypothesis enters.  For the record, the complete list of results
that the proof of Theorem~\ref{thm:main} uses as external (not re-proved here)
is: Theorem~\ref{thm:knmilnor} (Kneser--Milnor),
Theorem~\ref{thm:gl} (integrality of reducing slopes), Theorem~\ref{thm:propP}
(Property P), Theorem~\ref{thm:niwu} (the $L$-space cabling theorem), and
Theorem~\ref{thm:fibered} (fibredness of $L$-space knots, used only in the
$0$-surgery case), together with the sphere theorem and the classification of
closed orientable prime $3$-manifolds with finite cyclic fundamental group as
lens spaces.  Everything else---the homology computation, the primality of
nonintegral surgeries, the cabling-surgery splitting, and the whole case
analysis---is proved in Sections~\ref{sec:prelim}--\ref{sec:proof}.
\end{remark}

\end{document}
